\week{9}
% \section*{Week 9. Lecture 22.}
\lecture{22}{Classifications of singularities. OH NOO....}
% \emoji{grinning-face}
In this lecture we are going to see more examples of classification of singularities and will study the zeros of an analytic function and a connection between zeros and poles.

\begin{example} 18.4:} What type of singularity does $f(z) = \frac{e^z - 1}{z}$ have at 0?

Sol: The Laurent series of $f(z)$ about $z=0$ is given by:
\[ f(z) = \frac{1}{z} \left( \sum_{n=0}^{\infty} \frac{z^n}{n!} - 1 \right) = \frac{1}{z} \left( 1 + z + \frac{z^2}{2!} + \frac{z^3}{3!} + \dots - 1 \right) = 1 + \frac{z}{2!} + \frac{z^2}{3!} + \dots \]
The principal part of the Laurent series is $= 0 \implies$ it is a removable singularity. Define 
\[ \tilde{f}(z) = \begin{cases} \frac{e^z - 1}{z} & z \neq 0 \\ 1 & z = 0 \end{cases}\] 
$ \implies \tilde{f}$ is analytic for every $z \in \mathbb{C}$ (entire).
\end{example}

Alternatively: By Prop 16.1: $\lim_{z \to 0} z \cdot \frac{e^z - 1}{z} = \lim_{z \to 0} (e^z - 1) = 0 \implies z = 0$ is a removable singularity.

\begin{example} 18.5:} Find all the singularities and determine their nature:
1) $f(z) = \frac{z+1}{z^2 + 9}$
2) $f(z) = \frac{z^3 + 2z}{(z-i)^3}$

Sol: 1) Since $z^2 + 9 = (z+3i)(z-3i)$, $f(z)$ has singularities at $z = \pm 3i$. For $z = 3i$: Rewrite $f(z) = \frac{\varphi(z)}{z-3i}$, where $\varphi(z) = \frac{z+1}{z+3i}$ is analytic in a neighborhood of $z = 3i$ (and at $3i$), and $\varphi(3i) = \frac{1}{2} - i\frac{1}{6} \neq 0 \implies$ by Prop 16.2 $f(z)$ has a simple pole at $z = 3i$, and $\text{Res}(f; 3i) = \varphi(3i) = \frac{1}{2} - i\frac{1}{6}$.

In the same way $z = -3i$ is a simple pole and $\text{Res}(f; -3i) = \frac{1}{2} + i\frac{1}{6}$.

2) $f$ has an isolated singularity at $z = i$. Rewrite: $f(z) = \frac{\varphi(z)}{(z-i)^3}$, where $\varphi(z) = z^3 + 2z$ is analytic in a neighborhood of $z = i$ (including $z=i$) and $\varphi(i) = i^3 + 2i = i \neq 0$. Therefore, by Prop. 16.2 $f$ has a pole of order $m = 3$ at $z = i$, and by Prop 16.2 $\text{Res}(f; i) = \frac{\varphi^{(2)}(i)}{2!} = \frac{6i}{2} = 3i$.
\end{example}

\begin{example} 18.6:} Compute the residue of $f(z) = \frac{\sin z}{z^4}$.
    
Sol: $f$ has an isolated singularity at $z = 0$. We use Taylor series expansion for $\sin z$ (about $z = 0$) to obtain:
\begin{align*} 
    f(z) & = \frac{1}{z^4} \left( z - \frac{z^3}{3!} + \frac{z^5}{5!} - \dots \right) = \frac{1}{z^3} - \frac{1}{3!} \frac{1}{z} + \frac{z}{5!} - \dots = \\
            & = \frac{1}{z^3}\left(1 - \frac{z^2}{3!} + \frac{z^4}{5!} - \dots \right) 
        \end{align*}
$\implies z = 0$ is a pole of order $  m = 3$ not $4$!) 
From the Laurent series we see that $\text{Res}(f; 0) = b_1 = -\frac{1}{3!} = -\frac{1}{6}$.
\end{example}

The importance of residues will become much clearer a bit later, when we will investigate the complex integration, but note that the term $\frac{1}{z-z_0}$ is the tricky one to integrate! Now we are going to study zeros of the analytic functions and a bit more poles.

\section{Zeros and Poles}
\begin{definition} 16.3} A point $z = z_0$ is called a zero of the function $f(z)$ if $f(z_0) = 0$. If at $z = z_0$, $f(z_0) = f'(z_0) = \dots = f^{(n-1)}(z_0) = 0$ and $f^{(n)}(z_0) \neq 0$, then $z_0$ is called a zero of order $n$ of $f(z)$.

If $n = 1$, then the zero is called a simple zero.
\end{definition}

\begin{example} 18.7:} For the function $f(z) = (z-3)^2(z+2)^3$: $z = 3$ is a zero of order 2 and $z = -2$ is a zero of order 3. Let us check the case $z = 3$:
\begin{align*}
    f'(z) &= 2(z-3)(z+2)^3 + 3(z-3)^2(z+2)^2 \implies f'(3) = 0 \\
    f''(z) &= 2(z+2)^3 + 6(z-3)(z+2)^2 + 6(z-3)(z+2)^2 + 6(z-3)^2(z+2) \implies f''(3) = 250 \neq 0
\end{align*}
$\implies z = 3$ is a zero of order $n = 2$.
\end{example}

Rmk: To find all the zeros of $f(z)$ we need to solve the equation $f(z) = 0$.

\begin{example} 18.8:} The function $f(z) = e^z$ does not have any zeros.

Pf: We need to show that there are no solutions of $e^z = 0$. Let $z = x + iy$. By Euler's formula $e^z = e^{x+iy} = e^x(\cos y + i\sin y) \implies \begin{cases} e^x \cos y = 0 \\ e^x \sin y = 0 \end{cases}$ (*). Since $x \in \mathbb{R} \implies e^x \neq 0$. But $y \in \mathbb{R} \implies$ there is no solution to (*).
\end{example}

\textbf{Prop 16.3:} $z = z_0$ is a zero of order $n$ of $f(z)$ (an analytic function) $\iff f(z)$ can be represented as $f(z) = (z-z_0)^n \varphi(z)$, where $\varphi(z)$ is analytic at $z = z_0$ and $\varphi(z_0) \neq 0$.

Pf: [$\implies$] Assume $z = z_0$ is a zero of order $n$, namely $f(z_0) = f'(z_0) = \dots = f^{(n-1)}(z_0) = 0, f^{(n)}(z_0) \neq 0$
$\implies$ in the Taylor series expansion of $f(z)$ first $n$ terms vanish
and we have $f(z) = a_n(z-z_0)^n + a_{n+1}(z-z_0)^{n+1} + \dots = (z-z_0)^n (a_n + a_{n+1}(z-z_0) + \dots)$.

Denote: $\varphi(z) = a_n + a_{n+1}(z-z_0) + \dots$. Then, $f(z) = (z-z_0)^n \varphi(z)$, $\varphi(z)$ is analytic at $z_0$, $\varphi(z_0) = a_n \neq 0$ since $a_n = \frac{f^{(n)}(z_0)}{n!} \neq 0$.

[$\impliedby$] Assume $f(z) = (z-z_0)^n \varphi(z)$, $\varphi(z)$ is analytic at $z_0, \varphi(z_0) \neq 0$.

Expand $\varphi(z)$ into Taylor series about $z = z_0$: $f(z) = (z-z_0)^n (\tilde{a_0} + \tilde{a_1}(z-z_0) + \tilde{a_2}(z-z_0)^2 + \dots) = (z-z_0)^n \tilde{a_0} + \dots$.

From the last formula it is easy to see that $z = z_0$ is a zero of order $n$ of $f(z)$.

\begin{example} 18.9:} What is the order of the zero $z_0 = 0$ for:
    \begin{enumerate}
       \item  $f(z) = z(e^z - 1)$
       \item  $f(z) = \frac{z^{10}}{z - \sin z}$
    \end{enumerate}
Sol:
\begin{enumerate}
    \item We expand $e^z$ into Taylor series (about $z_0 = 0$) and get: 
    \begin{align*}
         z(1 + z + \frac{z^2}{2!} + \frac{z^3}{3!} + \frac{z^4}{4!} + \dots - 1) & = z^2 + \frac{z^3}{2!} + \frac{z^4}{3!} + \dots \\ 
        &   = z^2(1 + \frac{z}{2!} + \frac{z^2}{3!} + \dots) = z^2 \varphi(z).
    \end{align*}
Since $\varphi(0) = 1 \neq 0$ ($\varphi$ is analytic at $z_0 = 0$) by Prop 16.3 $z_0 = 0$ is a zero of order 2 of $f(z)$.
\item Plug the expansion of $\sin z$ (about $z_0 = 0$) into $f$ and get:
\begin{align*}
     \frac{z^{10}}{z - \sin z} & = \frac{z^{10}}{z - (z - \frac{z^3}{3!} + \frac{z^5}{5!} - \frac{z^7}{7!} + \dots)} \\ 
     & = \frac{z^{10}}{\frac{z^3}{3!} - \frac{z^5}{5!} + \frac{z^7}{7!} - \dots} = \frac{z^{10}}{z^3(\frac{1}{3!} - \frac{z^2}{5!} + \frac{z^4}{7!} - \dots)} 
\end{align*}
Factorize $z^3$.
Let 
\[ \varphi(z) = \frac{1}{3!} - \frac{z^2}{5!} + \frac{z^4}{7!} - \dots\]
Then $f(z) = z^7 \varphi(z)$, $\varphi$ is analytic at $z_0 = 0, \varphi(0) = \frac{1}{3!} = 6 \neq 0 \implies$ by Prop 16.3 $z_0 = 0$ is a zero of order 7 of $f(z)$.
\end{enumerate}
\end{example}

\begin{definition} 16.4:} $z = z_0$ is an isolated zero of $f(z)$ if there exists a neighborhood of $z_0$ such that $f(z)$ vanishes only at $z_0$ in this neighborhood.
\end{definition}
\textbf{Thm 16.1:} The zeros of an analytic function are isolated (all).

Pf: Assume $z = z_0$ is a zero of $f(z)$. Then, by Prop 16.3 $f(z) = (z-z_0)^n \varphi(z)$, where $\varphi(z)$ is analytic and $\varphi(z_0) \neq 0 \implies \varphi(z)$ is continuous at $z_0$ (Prop 6.5). Therefore, there exists a neighborhood of $z = z_0$ such that $\varphi(z) \neq 0$ for all $z$ in this neighborhood and in this neighborhood $f(z)$ has no other zeros different from $z = z_0$.

Next proposition provides a relation between zeros and poles of an analytic function.

\textbf{Prop 16.4:} If $p(z), q(z)$ are analytic at $z_0$ and $p(z_0) \neq 0$, then $f(z) = (\frac{p}{q})(z)$ has a pole of order $m$ at $z_0 \iff q$ has a zero of order $m$ at $z_0$.

Pf: [$\implies$] Suppose $f$ has a pole of order $m$ at $z_0$. Then, by Prop 16.2 $f(z) = \frac{\varphi(z)}{(z-z_0)^m}$, where $\varphi(z)$ is analytic at $z_0$ and $\varphi(z_0) \neq 0$. Thus, we have: $q(z) = (z-z_0)^m \frac{p(z)}{\varphi(z)}$, where $\frac{p(z_0)}{\varphi(z_0)} \neq 0$. Thus, $q$ has a zero of order $m$ at $z_0$ (by Prop 16.3).

[$\impliedby$] Assume $q$ has a zero of order $m$ at $z_0$. Then, by Prop 16.3 $q(z) = (z-z_0)^m \psi(z)$, where $\psi(z)$ is analytic at $z = z_0$ and $\psi(z_0) \neq 0$. Thus, $f(z) = \frac{p(z)}{(z-z_0)^m \psi(z)} = \frac{p(z)/\psi(z)}{(z-z_0)^m}$, where $\frac{p(z_0)}{\psi(z_0)} \neq 0$ ($\frac{p}{\psi}$ is analytic at $z = z_0$). Therefore, by Prop 16.2 $f$ has a pole of order $m$ at $z_0$.

\textbf{Cor 16.5:} If $p(z_0) \neq 0$ and $q(z)$ has a simple zero at $z_0$, then $\frac{p}{q}$ has a simple pole at $z_0$, and $\text{Res}(\frac{p}{q}, z_0) = \frac{p(z_0)}{q'(z_0)}$.

Pf: By Prop 16.4 with $m = 1$: $\frac{p}{q}$ has a simple pole at $z_0$.

To find the residue, we expand $q$ into Taylor series about $z_0$: $q(z) = q(z_0) + q'(z_0)(z-z_0) + \dots = q'(z_0)(z-z_0) + \frac{q''(z_0)}{2!}(z-z_0)^2 + \dots$, $q(z_0) = 0, q'(z_0) \neq 0$.

Let us rewrite the denominator and factor a common factor $z - z_0$ from the denominator:
\[  \frac{p(z)}{q(z)} = \frac{1}{z-z_0} \cdot \frac{p(z)}{q'(z_0) + \frac{q''(z_0)}{2!}(z-z_0) + \dots} \]
\begin{example} 18.10:} Compute all the residues of $f(z) = \frac{1}{z(e^z - 1)}$.

Sol: Observe that $f$ has infinitely many isolated singularities at the points $z = 2\pi ni, n \in \mathbb{Z}$. We shall break the argument into 2 cases:

$n = 0$: The singularity is at $z = 0$. From Ex 18.9 1) $z = 0$ is a double zero (of order 2) of $z(e^z - 1) \implies$ by Prop 16.4 $z = 0$ is a pole of order 2 of $f(z)$. 
The residue is the coefficient of $\frac{1}{z}$ in the 
Laurent series expansion of $f$: \[ f(z) = \frac{1}{z^2} \cdot \frac{1}{1 + \frac{z}{2!} + \frac{z^2}{3!} + \dots} \implies \varphi(z) = 1 + \frac{z}{2} + \frac{z^2}{3!} + \dots\] 
We get
\[ \text{Res}(f, 0) = \frac{\varphi'(0)}{1!} = -\frac{\frac{1}{2} + \frac{2z}{3!} + \frac{3z^2}{4!} + \dots}{\left(1 + \frac{z}{2} + \frac{z^2}{3!} + \dots \right)^2}\bigg|_{z=0} = -\frac{1}{2} \quad (m = 2)\]

\end{example}

$n \neq 0$: In terms of Prop 16.4 we have $p(z) = 1, q(z) = z(e^z - 1)$ and $z = 2\pi ni$ is a simple zero of $q(z)$. $q'(z) = ze^z + e^z - 1 \implies q'(2\pi ni) = 2\pi ni$ (Recall: $e^{2\pi ni} = 1$) $\implies f$ has a simple pole at $2\pi ni, n \neq 0$, and the residue is given by (Cor 16.5) $\frac{p(2\pi ni)}{q'(2\pi ni)} = \frac{1}{2\pi ni}$.

% \section*{Lectures 23 + 24}
\lecture{23+24}{Complex Integration. HUUURRAY>>>}
\section{Complex integration.}

In this lecture we first define the integral of a complex function of a real variable $\int w(t) dt$ and then the integral of a complex function of complex variable. Let us start with some motivation from real case.

Recall the definition of the integral of a real function: Let $f(x)$ be a real, continuous on a closed, bounded interval $[a, b]$. Partition the interval $[a, b]$ into $n$ subintervals: $a = x_0 < x_1 < \dots < x_{n-1} < x_n = b$, choose a sample point $x_j^*$, $0 \leq j \leq n$ in each subinterval. Then, $\int_a^b f(x) dx = \lim_{n \to \infty} \sum_{j=1}^n f(x_j^*)(x_j - x_{j-1})$, where the width of partition $\max |x_j - x_{j-1}| \to 0$.

Q-n: In the complex case, what will take the place of the partition?

In the complex plane there are a lot of curves that connect 2 points. So, the analogue of the real integration is:

For a given curve $C$ from $u + iv$ to $u' + iv'$ define $\int_C f(z) dz = \lim_{n \to \infty} \sum_{j=1}^n f(\xi_j)(z_j - z_{j-1})$, $\max |z_j - z_{j-1}| \to 0$, where we partition the curve and $(\xi_{j-1}, \xi_j)$ is an arc, $|z_j - z_{j-1}|$ is a length of a string that connects $z_{j-1}$ and $z_j$ on this arc.

If this limit is independent of the curve and of the choice of a point for the value $f(\xi_j)$, then this integral exists and it is called a path integral with respect to the path C. To see this first we need to know more about curves and paths.

\begin{definition} 17.1} A path is a continuous map $\gamma: [a, b] \to \mathbb{C}$ ($t \to \gamma(t)$), $\gamma(t) = x(t) + iy(t))$ where $[a, b] \subset \mathbb{R}$ is a closed interval. The image of any path $\gamma$ in $\mathbb{C}$ is called a curve. We say that $\gamma$ is a parametrization of that curve.
\end{definition}

\begin{example} 19.1} Define and compare 2 paths: $\gamma_1: [0, 1] \to \mathbb{C}$, $\gamma_1(t) = e^{2\pi i t}$; $\gamma_2: [0, 2\pi] \to \mathbb{C}$, $\gamma_2(t) = e^{it}$.

Observe: $\gamma_1(0) = \gamma_1(1) = \gamma_2(0) = \gamma_2(2\pi) = 1$. $\gamma_1, \gamma_2$ define the same curve as their images is the unit circle $\{z \in \mathbb{C}: |z| = 1\}$. But: $\gamma_1$ and $\gamma_2$ are different paths! $\gamma_1$ and $\gamma_2$ are parametrizations of the same curve!
\end{example}

\begin{example} 19.2:} Another useful example is the straight line connecting two points $w$ and $w'$ in the complex plane.

Straight line connecting $w$ and $w'$ (starting at $w$): $\gamma_3: [0, 1] \to \mathbb{C} \quad \gamma_3(t) = tw' + (1-t)w$.
\end{example}

\begin{definition} 17.2:} C is a simple curve if there exists a parametrization $\gamma$ of C which is injective (one to one): if $t_1 \neq t_2$, then $\gamma(t_1) \neq \gamma(t_2)$. C is a simple, closed curve if it is parametrized by $\gamma: [a, b] \to \mathbb{C}$ with $\gamma(a) = \gamma(b)$, but $\gamma(t_1) \neq \gamma(t_2)$ for all other $t_1, t_2 \in [a, b], t_1 \neq t_2$.
\end{definition}

Jordan's Thm: Every simple closed curve divides $\mathbb{C}$ into 2 regions - "inside" and "outside" the curve.

Notation: $-\gamma$ stands for the negative of the path $\gamma$ - the path which gives the same curve as $\gamma$, but traversed in the opposite direction.

\begin{example} 19.3} Write down $-\gamma$, where $\gamma$ is the straight line connecting $w$ to $w'$.

Sol: From Ex 19.2: $\gamma_3(t) = tw' + (1-t)w$ - straight line that connects $w$ to $w'$. Using this: $-\gamma_3: [0, 1] \to \mathbb{C} \quad -\gamma_3(t) = (1-t)w' + tw$ ($0 \to w', 1 \to w$).
\end{example}

\begin{example} 19.4:} Write in the form $z(t) = x(t) + iy(t)$ the following curve: $9x^2 + 4y^2 = 36 \iff (\frac{x}{2})^2 + (\frac{y}{3})^2 = 1$.

Sol: This is an equation of an ellipse with focal points $x = 2, y = 3$, since $\frac{x^2}{4} + \frac{y^2}{9} = 1 \iff (\frac{x}{2})^2 + (\frac{y}{3})^2 = 1$.
\end{example}

Parametrization: Take $x = 2\cos t, y = 3\sin t$ for $0 \leq t < 2\pi$. Then $(\frac{2\cos t}{2})^2 + (\frac{3\sin t}{3})^2 = \cos^2 t + \sin^2 t = 1 \implies z(t) = 2\cos t + i3\sin t$.

\begin{example} 19.5:} Write parametrization $z(t) = x(t) + iy(t)$ of the curve, described by $|z-i| + |z+i| = 4$.

Sol: $|z-i| + |z+i| = 4 \iff |x + i(y-1)| + |x + i(y+1)| = 4$
$\iff \sqrt{x^2 + (y-1)^2} + \sqrt{x^2 + (y+1)^2} = 4$
$\iff x^2 + (y-1)^2 = (x^2 + (y+1)^2 + 16 - 8\sqrt{x^2 + (y+1)^2})$
$\iff 8\sqrt{x^2 + (y+1)^2} = 16 + (y+1)^2 - (y-1)^2 = 16 + 4y$
$\iff 2\sqrt{x^2 + (y+1)^2} = 4 + y \iff 4x^2 + 4y^2 + 8y + 4 = 16 + y^2 + 8y$
$\iff 4x^2 + 3y^2 = 12 \iff \frac{x^2}{3} + \frac{y^2}{4} = 1 \iff (\frac{x}{\sqrt{3}})^2 + (\frac{y}{2})^2 = 1$

Again an ellipse equation. Take $x = \sqrt{3}\cos t, y = 2\sin t, 0 \leq t < 2\pi \implies$ parametrization $z(t) = \sqrt{3}\cos t + i2\sin t$.
\end{example}

We can add paths in the following manner:

\begin{definition} 17.3:} Given two paths $\gamma_1, \gamma_2$ defined by $\gamma_1: [a, b] \to \mathbb{C}, \gamma_2: [b, c] \to \mathbb{C}$ such that $\gamma_1(b) = \gamma_2(b)$ \{namely, the endpoint of the first path coincide with the starting point of the second path \} we define the sum: $\gamma_1 + \gamma_2: [a, c] \to \mathbb{C}$: $t \to \begin{cases} \gamma_1(t) \text{ if } t \in [a, b] \\ \gamma_2(t) \text{ if } t \in [b, c] \end{cases}$. We can also add curves by choosing appropriate parametrizations.
\end{definition}

\begin{definition} 17.4} Let C be a curve parametrized by the path $\gamma: [a, b] \to \mathbb{C}$, $\gamma(t) = x(t) + iy(t)$ is called differentiable if the real functions $x(t)$ and $y(t)$ are differentiable on $[a, b]$. If $\gamma$ is differentiable on $(a, b)$ and $\gamma'(t) \neq 0$ for $t \in (a, b)$, then C is called a smooth curve.
\end{definition}

\begin{definition} 17.5:} A contour is a piecewise-smooth curve: a finite union of smooth curves joined end-to-end. The length of a contour C is defined to be $L = \int_a^b |\gamma'(t)| dt$, where $\gamma: [a, b] \to \mathbb{C}$ is any smooth parametrization of C.
\end{definition}

\checkthis{} The length is independent of parametrization!

\begin{example} 19.6:} Define $\gamma: [0, 1] \to \mathbb{C}$ by $\gamma(t) = (1+i)t$. Graph this contour and calculate its length.

Sol: $L = \int_0^1 |\gamma'(t)| dt = \int_0^1 \sqrt{1^2 + 1^2} dt = \sqrt{2}$.
\end{example}

Now we study how to integrate a complex function of a real variable.

\begin{definition} 17.6} Given a complex function $w(t) = u(t) + iv(t)$ of a real variable $t$, we define:
\[ \int_a^b w(t) dt = \int_a^b u(t) dt + i \int_a^b v(t) dt \]
whenever the 2 integrals on the right side exist \{for example, if both $u$ and $v$ are continuous or piecewise continuous\}.
\end{definition}

From standard results of real analysis we have:

\textbf{Prop 17.1:} Let $w, w_1, w_2: [a, b] \to \mathbb{C}$ be complex functions. Then
\begin{enumerate}
    \item For any $c \in \mathbb{C}: \int_a^b cw(t) dt = c \int_a^b w(t) dt$
    \item $\int_a^b (w_1(t) + w_2(t)) dt = \int_a^b w_1(t) dt + \int_a^b w_2(t) dt$
    \item For any $\alpha \in \mathbb{C}: \int_a^b \alpha w(t) dt = \alpha \int_a^b w(t) dt$
    \item If $W(t)$ is an antiderivative of $w(t)$, then $\int_a^b w(t) dt = W(b) - W(a)$
\end{enumerate}

Note that the second and third parts of this proposition say that the integral defined is a linear map from the vector space of complex functions of a real variable to the complex numbers.

\begin{example} 19.7:} Compute $\int_0^{\pi/4} e^{it} dt$.

Sol: $e^{it}$ has antiderivative $\frac{1}{i}e^{it} = -ie^{it} \implies$ by Prop 17.1 4) $\int_0^{\pi/4} e^{it} dt = -ie^{it}|_0^{\pi/4} = -i[e^{i\frac{\pi}{4}} - 1] = -i[\frac{1}{\sqrt{2}} + i\frac{1}{\sqrt{2}} - 1] = \frac{1}{\sqrt{2}} + i(1 - \frac{1}{\sqrt{2}})$.
\end{example}

The following proposition is a very useful estimate for the size of an integral of a complex function of a real variable. The proof is omitted.

\textbf{Prop 17.2} $|\int_a^b w(t) dt| \leq \int_a^b |w(t)| dt$.

Our next subject is contour integrals, namely the integral of some function along some contour.

\section{Contour Integrals}
Let C be a contour and $f$ a continuous complex-valued function defined on C. Suppose that C is parametrized by the path $\gamma: [a, b] \to \mathbb{C}$.

\begin{definition} 17.7:} The contour integral of $f$ along C is defined as $\int_C f(z) dz = \int_a^b f(\gamma(t)) \gamma'(t) dt$.
\end{definition}

The basic idea is change of variables. Take $z = \gamma(t)$, then $dz = \frac{d\gamma}{dt} dt$, so the integral on the LHS should be equal to the one on the RHS. 
But, of course, to make it rigorous one needs to show that the expression on the RHS gives the same value as the definition with sums \[ \lim_{n \to \infty} \sum_{j=1}^n f(z_j)(z_j - z_{j-1}) (|z_j - z_{j-1}| \to 0).\] 
The definition of contour integral is much easier to use in practice. Let us see some examples.

\begin{example} 19.8:} Let C be a semicircle parametrized by the path $\gamma(\theta) = 2e^{i\theta}, -\frac{\pi}{2} \leq \theta \leq \frac{\pi}{2}$. Let $f(z) = \bar{z}$. Then:
\begin{align*}
     \int_C f(z) dz &= \int_{-\pi/2}^{\pi/2} f(\gamma(\theta)) \gamma'(\theta) d\theta = \int_{-\pi/2}^{\pi/2} \overline{2e^{i\theta}} 2ie^{i\theta} d\theta = \int_{-\pi/2}^{\pi/2} 4i d\theta \\ 
    & = 4i(\frac{\pi}{2} - (-\frac{\pi}{2})) = 4\pi i 
\end{align*}
\end{example}

\begin{example} 19.9:} Let C be the same contour as in Ex 19.8, $f(z) = z^2$.
    \begin{align*}
 \int_C f(z) dz & = \int_{-\pi/2}^{\pi/2} (2e^{i\theta})^2 2ie^{i\theta} d\theta = \int_{-\pi/2}^{\pi/2} 8i e^{3i\theta} d\theta \\ 
 & = \frac{8i}{3i} e^{3i\theta} |_{-\pi/2}^{\pi/2} = \frac{8}{3} [e^{i\frac{3\pi}{2}} - e^{-i\frac{3\pi}{2}}] = \frac{8}{3} 2i \sin(\frac{3\pi}{2}) = -i\frac{16}{3} 
    \end{align*}
\end{example}

\begin{example} 19.10:} Let $f(z) = z^2$, C - the straight line segment from $-2i$ to $2i \implies$ C is parametrized by the path $\gamma(t) = 2it + (-2i)(1-t) = -2i + 4it, 0 \leq t \leq 1$. Then:
\begin{align*} 
    \int_C f(z) dz & = \int_0^1 f(\gamma(t)) \gamma'(t) dt \\ 
    & = \int_0^1 (-2i + 4it)^2 4i dt = 4i(-4t - \frac{16t^3}{3} + \frac{16t^2}{2})|_0^1 \\ 
    & = 4i(-4 - \frac{16}{3} + 8) = -i\frac{16}{3} 
\end{align*}
\end{example}

Note: in the last two examples we have integrated $f(z) = z^2$ on two different contours that connect $2i$ to $-2i$ and obtained the same value, this is not a coincidence!

Prop 17.1 + basic linearity properties of integration applied to $\int_C f(z) dz$ imply:

\textbf{Prop 17.3:} Let $f, g$ be complex functions and $C, C_1, C_2$ contours.
\begin{enumerate}
    \item $\int_{C_1 + C_2} f(z) dz = \int_{C_1} f(z) dz + \int_{C_2} f(z) dz$ and $\int_{-C} f(z) dz = -\int_C f(z) dz$
    \item $\int_C (f+g)(z) dz = \int_C f(z) dz + \int_C g(z) dz$
    \item For any $\alpha \in \mathbb{C}: \int_C \alpha f(z) dz = \alpha \int_C f(z) dz$
    \item If $f$ is continuous on a domain $D \subseteq \mathbb{C}$ and has antiderivative $F(z)$ on D, then $\int_C f(z) dz = F(z_1) - F(z_0)$, where $z_0$ and $z_1$ are the endpoints of C.
\end{enumerate}

Pf: 1), 2), 3) follow directly from Prop 17.1.

4) If F is an antiderivative of $f$, then

\[ \frac{d}{dt} F(\gamma(t)) = F'(\gamma(t)) \gamma'(t) = f(\gamma(t)) \gamma'(t) \]
Here $\gamma: [a, b] \to \mathbb{C}$ is a parametrization of a contour C in D with $\gamma(a) = z_0$ and $\gamma(b) = z_1$. Applying Prop 17.1 with $w(t) = f(\gamma(t))\gamma'(t)$ and $F(\gamma(t)) = W(t)$ yields the result.

\textbf{Jordan Curve Theorem:} Let C be simple closed curve (Jordan curve) in $\mathbb{R}^2$. Then its complement $\mathbb{R}^2 \setminus C$ consists of exactly two connected components. One of these components is bounded (the interior) and the other is unbounded (the exterior), and the curve C is the boundary of each component.

The proof is highly non-trivial and it is omitted.
